\subsection{Nilpotent ideals with regular Rees quotients}

\begin{lemma}\label{E:regular-balanced-finiteness}
Let $B$ be a finite regular pointed monoid, with its zero as
basepoint, and put $R_B=\Sigma^\infty B$.  Let $p$ divide none of
the orders of its nonzero maximal subgroups.  If $L,V$ are a right
and a left $R_B$-module with finite-dimensional total
$\mathbb F_p$-homology, then
$L\wedge_{R_B}^{\mathbb L}V$ also has finite-dimensional total
$\mathbb F_p$-homology.
\end{lemma}

\begin{proof}
Induct on the number of nonzero $\mathcal J$-classes of $B$.
Choose a minimal nonzero class $J$ and an idempotent $e\in J$.
Since $B$ is regular, $J\cup\{0\}$ is a completely $0$-simple
principal ideal.  If $H$ is the maximal group at $e$, then
$eBe=H\cup\{0\}$ and multiplication gives a based bijection
\[
 Be\wedge_H eB\ \cong\ J\cup\{0\}.                 \tag{E.39}
\]
This is the Rees matrix normal form: after normalizing the sandwich
entry at $e$ to $1$, each element $(i,h,j)$ is uniquely a product
$(i,h,0)(0,1,j)$ modulo the intermediate $H$-coordinate.
The nonbasepoint parts of $Be$ and $eB$ are free on their indicated
$H$-sides.  Thus the derived balanced smash in (E.39) agrees with
the ordinary one after taking suspension spectra.

Write $R=R_B$, $K=\Sigma^\infty(J\cup\{0\})$, and
$A=\Sigma^\infty(B/J)$.  Then
\[
 K\simeq Re\wedge_{\mathbb S[H]}^{\mathbb L}eR,
 \qquad K\longrightarrow R\longrightarrow A,
 \qquad A\wedge_R^{\mathbb L}A\simeq A.               \tag{E.40}
\]
The last statement follows by tensoring the preceding cofiber
sequence, since $Ae=eA=0$.  Associativity gives the natural triangle
\[
 Le\wedge_{\mathbb S[H]}^{\mathbb L}eV
 \longrightarrow L\wedge_R^{\mathbb L}V
 \longrightarrow
 (L\wedge_R^{\mathbb L}A)
       \wedge_A^{\mathbb L}(A\wedge_R^{\mathbb L}V).   \tag{E.41}
\]
The idempotent summands $Le,eV$ have finite total mod-$p$ homology.
Extension to $H\mathbb F_p$ commutes with derived tensors and
homotopy orbits; the resulting modules have homology in a bounded
range, so the following group-homology spectral sequence converges.
The first term in (E.41) is the homotopy orbit spectrum of their
smash product under the finite group $H$.  The group-homology
spectral sequence collapses to coinvariants because $p\nmid|H|$;
its total mod-$p$ homology is therefore finite-dimensional.

The descended modules also retain finite total mod-$p$ homology.
For example, $L\wedge_R^{\mathbb L}A$ is the cofiber of
$Le\wedge_{\mathbb S[H]}^{\mathbb L}eR\to L$.
As a left $H$-set, $eB\setminus\{0\}$ is finite and free, so the
source is, on underlying spectra, a finite wedge of copies of $Le$.
The analogous assertion holds on the left.  The quotient $B/J$ is
again a finite regular pointed monoid with fewer nonzero classes,
and its nonzero maximal groups are among those of $B$.
Induction makes the last term in (E.41) finite in total mod-$p$
homology.  The long exact sequence proves the assertion for its
middle term.
\end{proof}

\begin{theorem}\label{E:nilpotent-ideal-regular-quotient}
Let $M$ be a finite monoid with rectangular-band minimal ideal $K$.
Suppose the pointed Rees quotient $Q=M/K$ has a nilpotent two-sided
ideal $N$ containing zero such that $Q/N$ is a regular pointed
monoid.  No section of $Q\to Q/N$ is assumed.
If $p$ divides none of the orders of the nonzero maximal subgroups
of $Q/N$, then $H^*(BM;\mathbb F_p)$ is locally finite over the
Steenrod algebra.  In particular $BM\not\simeq\Sigma\CP^\infty$.

For a minimal Rees matrix ideal over a finite group $H$, the same
conclusion holds provided also that $p\nmid|H|$.
\end{theorem}

\begin{proof}
First assume the minimal ideal is a band $I\times J$.  Use
Proposition~\ref{E:band-free-group-bar} to write
\[
 \widetilde{\Sigma^\infty BM}\simeq
 Z:=Y\wedge_A^{\mathbb L}X,
 \quad A=\Sigma^\infty Q,
 \quad X=\Sigma^\infty BF_I,
 \quad Y=\Sigma^\infty BF_J.                          \tag{E.42}
\]
We construct actual filtered ring and module spectra, rather than
only a filtration of their ordinary homology modules.

Choose $h\geq1$ such that $N^{h+1}=\{0\}$.  Give $A$ the finite
decreasing filtration
\[
 A^0=A,\qquad A^s=\Sigma^\infty N^s\quad(s\geq1),
 \qquad A^{h+1}=0.                                   \tag{E.43}
\]
Every $N^s$ is a pointed ideal, and products satisfy
$N^sN^t\subseteq N^{s+t}$.  Thus (E.43) is a filtered associative
ring spectrum, with the sphere unit in weight zero.
For the endpoints put
\[
 X^0=X,\qquad
 X^s=\Sigma^\infty
       \bigcup_{a\in N^s\setminus\{0\}}
              B\operatorname{im}\Phi_a\quad(s\geq1),  \tag{E.44}
\]
and define $Y^s$ in the same way for the right action.  An empty
union means the basepoint.  These are actual submodule spectra
under all of $Q$: composing any actor with an element of $N^s$
again gives an element of $N^s$, or zero.  More generally the
action maps factor strictly as
\[
 A^s\wedge X^t\longrightarrow X^{s+t},\qquad
 Y^t\wedge A^s\longrightarrow Y^{s+t},                 \tag{E.45}
\]
with the usual interpretation in weight zero.  These factorizations
respect all multiplication identities.  We therefore have coherent
filtered left and right modules, not merely maps whose homotopy
classes happen to satisfy (E.45).

All terms of these three filtrations are connective.  The endpoint
terms have finite CW homotopy type, by
Lemma~\ref{E:free-subset-intersections}; the ring terms are finite
wedges of spheres.  Their associated graded terms consequently
have finite-dimensional total mod-$p$ homology.

\emph{The filtered absolute bar.}
Take the ordinary derived bar over the sphere, with degree-$n$
term $Y\wedge A^{\wedge n}\wedge X$.  Use convolution of the
decreasing filtrations: the level-$L$ part is the homotopy colimit
of tensors of filtration terms whose indices sum to at least $L$.
For fixed $L$ this is a finite colimit calculation, since every
input filtration is zero above $h$; zero tails may equivalently
be truncated at $\max(h+1,L)$ in each index.
Equations (E.43)--(E.45) make all bar faces and degeneracies
filtered.  Let $F^LZ$ be the realization of its level-$L$ part.
Its level zero is $Z$ in (E.42).

In bar degree $n$ every nonzero filtration index is at most $h$,
and there are $n+2$ factors.  Thus the level-$L$ bar term is zero
whenever $L>(n+2)h$.  Every remaining term is connective, because
connective spectra are closed under smash products and colimits.
The skeletal realization filtration adds its degree-$n$ normalized
term with suspension $n$.  It follows that
\[
 F^LZ\text{ is at least }\lceil L/h\rceil-2
       \text{ connective}.                           \tag{E.46}
\]
In particular the connectivity tends to infinity with $L$.

\emph{Finite cohomology of each weight.}
For a decreasing filtration write
$\operatorname{gr}^s T=\operatorname{cofib}(T^{s+1}\to T^s)$.
The associated graded construction is strong monoidal for the
convolution just used:
\[
 \operatorname{gr}^w(U\wedge V)
 \simeq\bigvee_{s+t=w}
        \operatorname{gr}^sU\wedge\operatorname{gr}^tV. \tag{E.47}
\]
This follows by taking the successive cofibers of the finite
pushout diagrams defining convolution; smash products preserve
cofiber sequences.  Iteration gives the same formula for any
finite number of factors.  Associated graded also commutes with
geometric realization, since both cofibers and realization are
colimits.  Consequently $\operatorname{gr}Z$ is the bar over the
graded ring spectrum $\operatorname{gr}A$, with endpoints
$\operatorname{gr}Y,\operatorname{gr}X$.

The weight-zero ring is
\[
 R=\operatorname{gr}^0A
   =\operatorname{cofib}(\Sigma^\infty N\to\Sigma^\infty Q)
   \cong\Sigma^\infty(Q/N).                           \tag{E.48}
\]
In the \emph{graded} ring there is a canonical inclusion of this
weight-zero part.  Thus $\operatorname{gr}A=R\vee D$, where $D$
is supported in the strictly positive weights $1,\ldots,h$.
This splitting is in the associated graded ring; no multiplicative
section into $A$ is being asserted.

Compute the same graded derived tensor using the relative bar
over $R$.  Its normalized degree-$n$ term is
\[
 \operatorname{gr}Y\wedge_R^{\mathbb L}
      D\wedge_R^{\mathbb L}\cdots\wedge_R^{\mathbb L}D
      \wedge_R^{\mathbb L}\operatorname{gr}X,           \tag{E.49}
\]
with $n$ copies of $D$.  To see the normalization directly, expand
each internal copy of $R\vee D$.  Summands containing an $R$ are
exactly the degeneracy summands, and their latching cofiber is
the all-$D$ summand (E.49).

In a fixed weight $w$, (E.49) is a finite wedge of tensors
\[
 (\operatorname{gr}^{s_0}Y)\wedge_R^{\mathbb L}
 (\operatorname{gr}^{s_1}A)\wedge_R^{\mathbb L}\cdots
 \wedge_R^{\mathbb L}(\operatorname{gr}^{s_n}A)
 \wedge_R^{\mathbb L}(\operatorname{gr}^{s_{n+1}}X),
\]
where the indices sum to $w$, the two endpoint indices are
nonnegative, and all $n$ internal indices are positive.
Each factor has finite total mod-$p$ homology.  Since $Q/N$ is
regular, repeated application of
Lemma~\ref{E:regular-balanced-finiteness} gives the same finiteness
for each displayed tensor.  Moreover $n\leq w$, so only finitely
many normalized relative bar degrees occur in weight $w$.
Its finite skeletal filtration proves
\[
 \dim_{\mathbb F_p}H^*(\operatorname{gr}^wZ;\mathbb F_p)
       <\infty\qquad\text{for every }w\geq0.          \tag{E.50}
\]

For fixed $L$, the cofiber $Z/F^LZ$ has a finite filtration with
successive cofibers $\operatorname{gr}^wZ$, $0\leq w<L$.
Its total mod-$p$ cohomology is therefore finite-dimensional.
By (E.46), every class of $H^d(Z;\mathbb F_p)$, for fixed $d$,
is pulled back from one such cofiber when $L$ is sufficiently
large.  The cofiber map is an actual map of spectra.  Stable
Steenrod operations commute with it, so the entire orbit of the
class lies in the image of the finite-dimensional total
cohomology of this one cofiber.  This proves local finiteness.
The associated graded argument was used only for finite dimension
estimates; no identification of Steenrod operations on an
algebraic spectral sequence is required.

Equation (E.42) transfers the conclusion to $BM$.  If the minimal
ideal has a nontrivial group $H$, its actual classwise-averaging
quotient is a mod-$p$ homology equivalence for $p\nmid|H|$, as
proved in Corollary~\ref{E:two-regular-classes}.  It leaves $Q=M/K$
unchanged, so the band case applies.  Finally, a prime avoids all
the finitely many specified group orders, while the suspended
generator of $H^2(\CP^\infty;\mathbb F_p)$ has an infinite
Steenrod orbit at every prime.
\end{proof}

\begin{remark}\label{E:regular-quotient-scope}
This theorem includes the split regular-submonoid situation
$Q=B\cup N$, $B\cap N=\{0\}$, as well as the earlier finite-unit
case.  Its hypothesis still restricts the ideal structure.
A nonregular element may have a nonzero regular element below it
in the $\mathcal J$-order, preventing the nonregular elements from
lying in a nilpotent ideal.  Removing additional regular ideals
preserves finite underlying endpoint spectra, but does not by
itself supply the finite image filtrations (E.44) for their
coherent quotient actions.  The principal-ideal homotopy-colimit
construction of Lemma~\ref{E:regular-controlled-localization}
provides that additional step.  The resulting general theorem and
the finite-category reduction are proved in the following subsections.
\end{remark}
